% ------------------------------------------------------------
% Page 97
% ------------------------------------------------------------

\section*{L'Instituteur, le Maire et le Pasteur}

\subsection*{Mon arrivée aux Oubrets en 1938}

Je suis élève-maître sortant, et, avec tous mes camarades de promotion, je participe au dernier mouvement des instituteurs, sachant parfaitement que nous serons attribués des postes dont personne n'a voulu ! J'ai inscrit sur ma fiche des postes dans le sud du département.

Dans les premiers jours de ce mois de septembre, politiquement très secoué\footnote{Accords entre la Grande Bretagne, la France et Hitler, en septembre 1938, acceptant l'annexion des Sudètes par l'Allemagne (région de Tchécoslovaquie où les Allemands étaient une forte minorité). La paix était sauvée pour quelques mois., mais Churchill 1\textsuperscript{er} ministre anglais, au retour des plénipotentiaires alliés, n'hésita pas à dire « Nous aurons la guerre et le déshonneur d'avoir abandonné nos alliés. »}, je reçois ma nomination à Meyrueis (Les Oubrets), Meyrueis désignant le chef-lieu de commune et « Les Oubrets » le hameau. Je consulte la carte de la Lozère, et ne trouve pas Les Oubrets. C'est un peu inquiétant, mais Meyrueis me rassure : c'est même un canton ! A priori je suis satisfait et confiant ; on verra bien... !

Ce 28 Septembre donc, avec une valise lourdement chargée, je débarque à Meyrueis, vers 6 heures du soir. Le temps est maussade, il bruine. Je demande aussitôt la direction des Oubrets. On me renseigne : « Vous n'y êtes pas encore ! c'est à 12 km, aux sources de la Brèze, sous le col de la Serreyrède ». Ces 12 km me tombent sur la nuque ; que faire ? Dormir à Meyrueis, c'est retarder le problème. Je me décide pour un taxi.

Je trouve, difficilement, un taxi qui accepte de me conduire aux Oubrets. Aux réticences du chauffeur je comprends que la route ne doit pas être très bonne ; « Ce sera 45 francs » me précise-t-il. Il fait déjà nuit.

Nous voilà donc sur la petite route qui longe une rivière, la Brèze. Mon chauffeur n'est pas très bavard. Il m'indique toutefois que nous traverserons Pourcarès, à 4 km de Meyrueis. Je distingue au bord de la route une école qui aurait bien fait mon affaire.

Le goudron s'arrête là, et la route devient pentue, cahotante. Le chauffeur, de moins en moins bavard, me désigne le petit hameau de Campis. Très préoccupé par sa conduite, il répond à peine à mes questions, mais se concentre contre les ornières que, soucieux de sa mécanique, il s'efforce de chevaucher. « Il nous reste 5 km à faire, les plus mauvais », me dit-il, « et ça grimpe ! ».

\subsection*{L'accueil}

Il me dépose enfin devant une assez grande maison qui pourrait être l'école ; la maison d'en face est éclairée. Je frappe et je suis accueilli par deux jeunes hommes ; je n'ai pas le temps de me présenter ; « Vous êtes le nouvel instituteur, on vous attendait... comme tous les ans... on a l'habitude ! »
